% Generated draft: requires mathematical and source review.
\documentclass[11pt]{article}
\usepackage[margin=25mm]{geometry}
\usepackage{fontspec}
\setmainfont{Latin Modern Roman}
\usepackage{amsmath,amssymb,mathtools,unicode-math}
\setmathfont{Latin Modern Math}
\usepackage{enumitem,needspace,xurl,hyperref}
\setlength{\parindent}{0pt}
\setlength{\parskip}{5pt}
\setlength{\emergencystretch}{3em}
\newcommand{\N}{\mathbb{N}}
\newcommand{\Z}{\mathbb{Z}}
\newcommand{\Q}{\mathbb{Q}}
\newcommand{\R}{\mathbb{R}}
\newcommand{\C}{\mathbb{C}}
\newcommand{\F}{\mathbb{F}}
\newcommand{\A}{\mathbb{A}}
\newcommand{\PP}{\mathbb{P}}
\newcommand{\E}{\mathbb{E}}
\newcommand{\eps}{\varepsilon}
\newcommand{\dd}{\,\mathrm{d}}
\begin{document}
\section*{UnsolvedMath ID: 30004645}
\textbf{Source problem code:} catalogue entry
\subsection{Definitions and mathematical statement}
Let \(k\) be a field, let \(K/k\) be a finite field extension, and let \(x_1,\ldots,x_n\) be algebraically independent over \(K\). Write \(L=K(x_1,\ldots,x_n)\). Let a group \(G\) act on \(L\) by \(k\)-automorphisms, meaning that every group element fixes \(k\) pointwise.

The monomial-form interpretation of the supplied term \emph{quasi-monomial} is that every \(\sigma\in G\) preserves \(K\) and satisfies
\[
\sigma(x_j)=c_j(\sigma)\prod_{i=1}^{n}x_i^{a_{ij}(\sigma)}\qquad(1\leq j\leq n),
\]
where \(c_j(\sigma)\in K^{\times}\) and \(A(\sigma)=(a_{ij}(\sigma))\in\operatorname{GL}_n(\Z)\). The maps are required to constitute a group action. The source does not supply a definition of quasi-monomial action; in particular, whether its intended convention additionally requires \(K^G=k\) remains unresolved. Finiteness of \(G\), faithfulness of its action, and separability or normality of \(K/k\) are not explicitly stipulated and are not imposed here.

Define the fixed field by
\[
L^G=\{f\in L:\sigma(f)=f\text{ for every }\sigma\in G\}.
\]
A field extension \(F/k\) is \emph{\(k\)-rational} if there are an integer \(m\geq0\) and elements \(t_1,\ldots,t_m\in F\), algebraically independent over \(k\), such that \(F=k(t_1,\ldots,t_m)\); for \(m=0\), this means \(F=k\).

\textbf{Problem.} Determine necessary and sufficient conditions on the extension \(K/k\) and the specified quasi-monomial action of \(G\) for
\[
K(x_1,\ldots,x_n)^G
\]
to be \(k\)-rational. This is a request for a characterization, not an assertion of rationality for every such action. Its intended scope requires clarification of the quasi-monomial convention and the group hypotheses identified above.
\subsection{Short English statement}
Characterize when the field of rational functions fixed by a quasi-monomial group action is a purely transcendental extension of \(k\). The source leaves some action hypotheses implicit, particularly finiteness of \(G\) and whether \(K^G=k\) is required. The supplied open status is dated curation, not an independently verified conclusion.
\subsection{Sources}
\begin{itemize}
\item Supplied local UnsolvedMath registry record, permanent ID \(30004645\), title \emph{Rationality of Quasi-Monomial Fixed Fields}, source problem code \texttt{catalogue entry}, in the Oberwolfach Reports Open Problems collection. The record gives proposed year 2021 and names no proposer.
\item Supplied report citation: \emph{Homotopic and Geometric Galois Theory (2022)}, Oberwolfach Reports, pp.~663--744, DOI \texttt{10.4171/owr/2021/12}. \url{https://doi.org/10.4171/owr/2021/12}. The literature-source label is \emph{Oberwolfach Report 12/2021}. The supplied 2022 title date and 2021 report identification are retained without resolving their relationship.
\item Supplied literature reference: \emph{Rationality problem for quasi-monomial actions}. \url{https://repository.kulib.kyoto-u.ac.jp/dspace/bitstream/2433/232901/1/B51-14.pdf}. The record supplies no authors, publication date, or precise passage for this reference.
\item Literature assessment dated exactly \texttt{2026-08-22}: the supplied curation labels the problem open and states that the quasi-monomial rationality question contains Noether's problem and rationality of algebraic tori, that many cases are classified, and that no general necessary and sufficient criterion is known. This is supplied dated curation derived from the OWR record, not independently verified literature assessment; its status and cited claims require independent verification.
\item The registry labels the statement \texttt{corrected\_verified} and reports manual semantic, typographic, and LaTeX review. These are supplied editorial claims; this draft does not independently verify them or establish mathematical correctness.
\end{itemize}
\end{document}
